
This study proposes an automated framework for diagnosing failure types in large language models (LLMs) and routing failures to targeted correction strategies. Using attention entropy extracted from transformer models, the framework classifies entity-substitution errors versus other failure modes without human annotation. Experiments on TruthfulQA benchmark failures using GPT-2 demonstrate that entity-substitution errors exhibit significantly lower attention entropy over entity spans (mean = 0.062) compared to non-entity errors (mean = 0.300), yielding a statistically robust diagnostic signal (p < 0.001, Cohen's d = -1.13). Threshold-based classification achieves 86.7\% accuracy on held-out test data. The study validates measurement assumptions (NER F1 = 0.96, Wikipedia coverage = 1.0) and confirms classification utility using 73 processed benchmark failures. A synthetic correction experiment demonstrates structural feasibility of matched routing (entity-error to retrieval-augmented generation), achieving +24 percentage points over mismatched routing in configurable mock settings. Real-world correction effectiveness remains unvalidated; Wikipedia API-based retrieval and GPT-judge evaluation were not deployed. Limitations include single-model validation (GPT-2 only), manual gold labeling bottleneck (N=100), and 27\% sample loss due to tokenizer mismatch. The framework demonstrates that attention patterns can serve as automated diagnostic signals for benchmark-scale failure analysis.
